\chapter{Impetigo}

\section*{Definition}
Impetigo is a highly contagious superficial bacterial skin infection most common in children aged 2--5 years, presenting in two forms: non-bullous (honey-colored crusts) and bullous (intact, flaccid bullae). It is the most common bacterial skin infection in children worldwide.

\section*{Pathophysiology}
Impetigo begins with bacterial inoculation through minor skin breaks where organisms proliferate in the superficial epidermis. Non-bullous impetigo is caused by \textit{S. pyogenes} and/or \textit{S. aureus}; crusts are dried serum and neutrophils. Bullous impetigo is caused exclusively by \textit{S. aureus} producing exfoliative toxins A and B that cleave desmoglein-1 in the granular layer, causing intraepidermal cleavage and bulla formation.

\section*{Causes}
\begin{itemize}
  \item \textbf{Bacterial:} \textit{Staphylococcus aureus} (most common), \textit{Streptococcus pyogenes} (group A strep), or both
  \item \textbf{Risk factors:} Warm, humid climate, poor hygiene, crowded living conditions, pre-existing skin conditions (scabies, eczema, insect bites), participation in contact sports
  \item \textbf{Transmission:} Direct person-to-person contact, fomites, autoinoculation by scratching
\end{itemize}

\section*{When to See a Doctor}
See your doctor if:
\begin{itemize}
  \item Widespread lesions, multiple family members affected, or rapidly spreading
  \item Fever, malaise, or lymphadenopathy suggesting systemic involvement
  \item Dark red or purpuric lesions, pain out of proportion (suspect ecthyma or deep infection)
  \item Signs of post-streptococcal glomerulonephritis: hematuria, edema, hypertension
\end{itemize}

\section*{Related Symptoms}
\begin{itemize}
  \item Ecthyma (deep ulcerated pyoderma), folliculitis, cellulitis, lymphangitis, scarlet fever
\end{itemize}
\textbf{Important:} Post-streptococcal glomerulonephritis may follow impetigo by 2--3 weeks (particularly with nephritogenic M types 49 and 55), but prompt antibiotic treatment does not reliably prevent this complication.

\section*{Self-Care / Home Management}
\begin{itemize}
  \item Gently remove crusts with warm saline soaks followed by application of topical antibiotic
  \item Keep affected areas covered with gauze or clothing to prevent autoinoculation and spread to others
  \item Frequent hand washing; avoid sharing towels, bedding, and clothing
\end{itemize}

\section*{How Your Doctor Will Evaluate You}
\begin{itemize}
  \item Clinical diagnosis based on characteristic honey-colored crusts or flaccid bullae
  \item Gram stain and bacterial culture of exudate or bullous fluid for antibiotic guidance
  \item Urinalysis 2--3 weeks after infection if streptococcal etiology suspected (to screen for glomerulonephritis)
\end{itemize}

\section*{Treatment Options}
\begin{itemize}
  \item Mild, localized: Topical mupirocin 2\% ointment TID or retapamulin 1\% ointment BID for 5--7 days
  \item Widespread or bullous: Oral cephalexin 25--50 mg/kg/day in 4 divided doses, cefadroxil, or amoxicillin--clavulanate for 7 days
  \item MRSA risk: Clindamycin or trimethoprim--sulfamethoxazole based on local susceptibility patterns
  \item Children with impetigo should not attend school or daycare until 24--48 hours after initiating effective treatment
\end{itemize}

\clearpage
